\documentclass[12pt,a4paper]{report}

% ══════════════════════════════════════════════════════════════════════════
% PACKAGES
% ══════════════════════════════════════════════════════════════════════════

\usepackage[utf8]{inputenc}
\usepackage[T1]{fontenc}
\usepackage[french]{babel}
\usepackage[margin=2.5cm]{geometry}
\usepackage{mathptmx}              % Police proche de Times New Roman
\usepackage{setspace}              % Gestion de l'interligne
\usepackage{titlesec}              % Contrôle taille/style des titres
\usepackage{graphicx}
\usepackage{caption}
\usepackage{xcolor}
\usepackage{fancyhdr}
\usepackage{hyperref}
\usepackage{indentfirst}           % Garde la tabulation en début de paragraphe
\usepackage{tocloft}
\usepackage{array}
\usepackage{tikz}
\usepackage{pgfgantt}
\usetikzlibrary{shapes.geometric, arrows.meta, positioning, fit, backgrounds}

% ── Interligne 1.5 et texte justifié (par défaut) ────────────────────────────
\onehalfspacing

% ── Tailles des titres selon les consignes (14 / 14 / 13) ───────────────────
\titleformat{\section}{\normalfont\fontsize{14}{16}\bfseries}{\thesection}{1em}{}
\titleformat{\subsection}{\normalfont\fontsize{13}{15}\bfseries}{\thesubsection}{1em}{}
\titleformat{\subsubsection}{\normalfont\fontsize{12}{14}\bfseries}{\thesubsubsection}{1em}{}

% ── Légendes des figures/tableaux : taille 11 + italique ─────────────────────
\captionsetup{font={it,footnotesize},labelfont={it,footnotesize}}

% ── En-tête / pied de page ───────────────────────────────────────────────────
\pagestyle{fancy}
\fancyhf{}
\renewcommand{\headrulewidth}{0.4pt}
\renewcommand{\footrulewidth}{0.4pt}
\fancyfoot[L]{EL FELLAH Adil}
\fancyfoot[R]{Page \thepage}
\fancyhead[L]{\nouppercase{\leftmark}}

\hypersetup{
    colorlinks=true,
    linkcolor=black,
    citecolor=black,
    urlcolor=blue,
    pdftitle={Conception d'une architecture réseau sécurisée et plateforme e-services hôteliers},
    pdfauthor={EL FELLAH Adil}
}

% ══════════════════════════════════════════════════════════════════════════
% DOCUMENT
% ══════════════════════════════════════════════════════════════════════════

\begin{document}

% ─────────────────────────────────────────────────────────────────────────
% PAGE DE GARDE (pas de numérotation visible)
% ─────────────────────────────────────────────────────────────────────────

\begin{titlepage}
    \thispagestyle{empty}
    \begin{center}

        % ── Logos : à remplacer par les vrais fichiers logo-emsi.png ──
        \begin{minipage}{0.45\textwidth}
            \raggedright
            % \includegraphics[width=3cm]{logo-entreprise.png} % si applicable
        \end{minipage}
        \hfill
        \begin{minipage}{0.45\textwidth}
            \raggedleft
            % \includegraphics[width=3cm]{logo-emsi.png}
            \textbf{ÉCOLE MAROCAINE DES SCIENCES DE L'INGÉNIEUR}\\
            \small Membre de HONORIS UNITED UNIVERSITIES
        \end{minipage}

        \vspace{1cm}

        \rule{\textwidth}{1pt}

        \vspace{0.5cm}

        {\Large \textbf{INGÉNIERIE INFORMATIQUE ET RÉSEAUX}}\\

        \vspace{0.3cm}

        \textbf{ANNÉE UNIVERSITAIRE}\\
        \textbf{2025-2026}

        \vspace{1.5cm}

        {\huge \textbf{Projet de Fin d'Année}}

        \vspace{1cm}

        \textbf{\Large THÈME}

        \vspace{0.3cm}

        \fbox{
            \begin{minipage}{0.85\textwidth}
                \centering
                \vspace{0.3cm}
                {\Large \textbf{Conception d'une architecture réseau sécurisée par chambre
                et développement d'une plateforme d'e-services hôteliers
                (Réservation, Réclamations et Room Service)}}
                \vspace{0.3cm}
            \end{minipage}
        }

        \vspace{2cm}

        \begin{minipage}{0.45\textwidth}
            \textbf{RÉALISÉ PAR :}\\
            EL FELLAH Adil
        \end{minipage}
        \hfill
        \begin{minipage}{0.45\textwidth}
            \textbf{ENCADRÉ PAR :}\\
            M./Mme XXXXXXX (EMSI)
        \end{minipage}

        \vspace{1cm}

        \textbf{SOUTENU LE :} XX/XX/2026

    \end{center}
\end{titlepage}

% ─────────────────────────────────────────────────────────────────────────
% NUMÉROTATION ROMAINE (dédicace → sommaire)
% ─────────────────────────────────────────────────────────────────────────

\pagenumbering{roman}
\setcounter{page}{2}

% ── Dédicace ──────────────────────────────────────────────────────────────

\chapter*{Dédicace}
\addcontentsline{toc}{chapter}{Dédicace}

À mes parents, dont le soutien indéfectible a rendu possible chaque étape
de ce parcours. À mes professeurs de l'École Marocaine des Sciences de
l'Ingénieur, pour la qualité de l'enseignement dispensé. À tous ceux qui
ont contribué, de près ou de loin, à l'aboutissement de ce travail.

\clearpage

% ── Remerciements ─────────────────────────────────────────────────────────

\chapter*{Remerciements}
\addcontentsline{toc}{chapter}{Remerciements}

Je tiens à exprimer mes sincères remerciements à mon encadrant pédagogique
pour ses orientations avisées et sa disponibilité tout au long de ce projet.
Mes remerciements s'adressent également à l'ensemble du corps enseignant du
département Ingénierie Informatique et Réseaux de l'EMSI, ainsi qu'à toute
personne ayant contribué à l'enrichissement de ce travail par ses remarques
et ses suggestions.

\clearpage

% ── Résumé ────────────────────────────────────────────────────────────────

\chapter*{Résumé}
\addcontentsline{toc}{chapter}{Résumé}

Ce projet de fin d'année porte sur la conception d'une architecture réseau
sécurisée par chambre et le développement d'une plateforme web d'e-services
hôteliers. La partie réseau consiste à segmenter le réseau de l'hôtel par
VLAN (Virtual Local Area Network) afin d'isoler le trafic de chaque chambre
et de sécuriser les communications inter-services. La partie applicative
repose sur le framework Laravel 12 et propose quatre espaces utilisateurs
distincts : client, serveur, technicien et administrateur. La plateforme
prend en charge la réservation de chambres et de services avec vérification
de disponibilité en temps réel, la gestion du room service, ainsi que le
traitement des réclamations techniques avec priorisation automatique.

\vspace{0.5cm}

\noindent\textbf{Mots-clés :} architecture réseau, VLAN, sécurité, hôtellerie,
plateforme web, Laravel, e-services.

\clearpage

% ── Abstract ──────────────────────────────────────────────────────────────

\chapter*{Abstract}
\addcontentsline{toc}{chapter}{Abstract}

This final-year project covers the design of a per-room secure network
architecture and the development of a hotel e-services web platform. The
network part involves segmenting the hotel infrastructure using VLANs to
isolate each room's traffic and enforce inter-service security policies. The
application layer is built on the Laravel 12 framework and provides four
distinct user spaces: guest, waiter, technician, and administrator. The
platform supports room and service bookings with real-time availability
checking, room-service order management, and technical complaint handling
with automatic priority assignment.

\vspace{0.5cm}

\noindent\textbf{Keywords:} network architecture, VLAN, security, hospitality,
web platform, Laravel, e-services.

\clearpage

% ── Liste des abréviations ────────────────────────────────────────────────

\chapter*{Liste des abréviations}
\addcontentsline{toc}{chapter}{Liste des abréviations}

\begin{description}
    \item[BDD :] Base de données
    \item[EMSI :] École Marocaine des Sciences de l'Ingénieur
    \item[HTTP :] HyperText Transfer Protocol
    \item[MVC :] Modèle-Vue-Contrôleur
    \item[PFA :] Projet de Fin d'Année
    \item[SGBD :] Système de Gestion de Base de Données
    \item[SQL :] Structured Query Language
    \item[VLAN :] Virtual Local Area Network
    \item[VLSM :] Variable Length Subnet Masking
    \item[WI-FI :] Wireless Fidelity
\end{description}

\clearpage

% ── Liste des figures ─────────────────────────────────────────────────────

\listoffigures
\addcontentsline{toc}{chapter}{Liste des figures}

\clearpage

% ── Liste des tableaux ────────────────────────────────────────────────────

\listoftables
\addcontentsline{toc}{chapter}{Liste des tableaux}

\clearpage

% ── Table des matières ────────────────────────────────────────────────────

\tableofcontents

\clearpage

% ─────────────────────────────────────────────────────────────────────────
% NUMÉROTATION ARABE (à partir de l'introduction générale)
% ─────────────────────────────────────────────────────────────────────────

\pagenumbering{arabic}
\setcounter{page}{1}

% ── Introduction générale ─────────────────────────────────────────────────

\chapter*{Introduction générale}
\addcontentsline{toc}{chapter}{Introduction générale}

Le secteur hôtelier connaît une transformation numérique profonde, dans
laquelle la qualité de l'infrastructure réseau et la disponibilité de
services numériques constituent des facteurs différenciateurs majeurs. Dans
ce contexte, l'hôtel Rambla Marrakech a servi de cadre à un projet visant
à répondre simultanément à deux problématiques complémentaires : la
sécurisation des communications réseau par chambre et la dématérialisation
des services offerts aux clients. La plateforme développée centralise la
réservation de chambres et de prestations, la commande de room service, la
remontée de réclamations techniques ainsi que leur traitement par le
personnel de l'établissement.

Ce rapport est organisé en quatre chapitres :

\begin{itemize}
    \item Le premier chapitre présente le contexte général du projet, la
    problématique ainsi que le cahier des charges.

    \item Le deuxième chapitre est consacré à la conception de l'architecture
    réseau sécurisée par chambre.

    \item Le troisième chapitre détaille la conception et le développement de
    la plateforme e-services hôteliers.

    \item Le quatrième chapitre présente les tests réalisés, les résultats
    obtenus ainsi que les perspectives d'amélioration.
\end{itemize}

\clearpage

% ═════════════════════════════════════════════════════════════════════════
% CHAPITRE 1 — CONTEXTE GÉNÉRAL ET CADRE DU PROJET
% ═════════════════════════════════════════════════════════════════════════

\chapter{Contexte général et cadre du projet}

\section{Introduction}

Ce premier chapitre présente le cadre général dans lequel s'inscrit le
projet. Il expose la problématique identifiée, les objectifs poursuivis,
le cahier des charges fonctionnel et les choix méthodologiques et
technologiques qui ont guidé la conception et le développement.

\section{Contexte du projet}

L'hôtel Rambla Marrakech est un établissement de standing situé au c\oe ur
de la ville ocre. Comme de nombreux hôtels de taille intermédiaire, il fait
face à deux enjeux structurels : d'une part, l'absence d'une infrastructure
réseau segmentée exposant les équipements des chambres à des risques de
sécurité mutuelle ; d'autre part, la gestion encore manuelle ou téléphonique
de services tels que la réservation, les commandes au restaurant et les
demandes de maintenance, engendrant des délais de traitement et un manque
de traçabilité. Ce projet répond à ces deux enjeux de manière conjointe, en
proposant une solution réseau et applicative intégrée.

\subsection{Problématique}

La problématique centrale peut se formuler ainsi : comment concevoir, pour
un hôtel de taille intermédiaire, une infrastructure réseau qui garantisse
l'isolation du trafic de chaque chambre, tout en déployant une plateforme
web permettant à chaque profil d'utilisateur --- client, personnel de
restaurant, technicien et administrateur --- d'accéder uniquement aux
fonctionnalités qui lui sont dédiées, le tout de façon sécurisée et
opérationnelle ?

\subsection{Objectifs du projet}

Les objectifs du projet se déclinent en deux volets. Sur le plan réseau,
il s'agit de concevoir un plan d'adressage IP avec segmentation VLAN par
chambre et par étage, de définir les règles de filtrage inter-VLAN par
listes de contrôle d'accès (ACL), et de modéliser la topologie sous Cisco
Packet Tracer. Sur le plan applicatif, les objectifs sont : développer une
plateforme web multi-rôles sous Laravel 12 permettant la réservation de
chambres et de services hôteliers, la gestion du room service avec suivi
des commandes, la remontée et le traitement des réclamations techniques,
et l'authentification sécurisée incluant une connexion via OAuth Google.

\section{Cahier des charges}

Le cahier des charges fonctionnel regroupe les exigences suivantes.

\textbf{Exigences réseau :} segmentation VLAN par chambre (un VLAN par
étage au minimum) ; isolation du trafic client vis-à-vis du réseau
d'administration ; plan d'adressage IP structuré avec VLSM ; règles ACL
empêchant les accès non autorisés entre VLANs ; authentification Wi-Fi
différenciée par VLAN ; modélisation complète sous Cisco Packet Tracer.

\textbf{Exigences applicatives :} authentification des utilisateurs par
formulaire et par OAuth Google (via Laravel Socialite) ; contrôle d'accès
par rôle (client, serveur, technicien, admin) appliqué via middleware ;
module de réservation de chambres et de services avec vérification de
disponibilité en temps réel et détection de chevauchements de dates ;
module de room service avec carte de menu groupée par catégorie, création
de commandes et suivi d'état ; module de réclamations techniques avec
catégorisation, priorisation automatique et assignation à un technicien ;
interface d'administration des réservations et des réclamations ; galerie
photographique des chambres avec filtrage de disponibilité par dates.

\textbf{Exigences non fonctionnelles :} interface responsive cohérente avec
un système de design custom (tokens de couleurs, typographie Fraunces et
Public Sans) ; protection CSRF sur tous les formulaires ; pagination et
séparation stricte des espaces par middleware d'authentification.

\section{Méthodologie de gestion de projet (Scrum)}

\subsection{Présentation du framework Scrum}

La réalisation de la plateforme Rambla Marrakech a suivi le framework agile
\textbf{Scrum}, dont le principe central consiste à découper le développement
en itérations courtes et planifiées appelées \textit{sprints}. Chaque sprint
produit un incrément fonctionnel et testable, permettant un ajustement
continu des priorités en fonction des retours obtenus.

Les rôles définis dans ce cadre sont les suivants :

\begin{itemize}
    \item \textbf{Product Owner} : responsable de la définition et de la
    priorisation du \textit{backlog} produit, garant de la valeur fonctionnelle
    de chaque livraison.
    \item \textbf{Scrum Master} : facilitateur chargé de lever les obstacles
    et de veiller au respect des cérémonies Scrum (planning, revue, rétrospective).
    \item \textbf{Équipe de développement} : réalise les éléments du backlog
    sélectionnés pour chaque sprint.
\end{itemize}

\subsection{Déroulement des sprints}

Le développement a été organisé en cinq sprints successifs, chacun ciblant
un module fonctionnel précis :

\begin{itemize}
    \item \textbf{Sprint 0 (Initialisation)} : mise en place de l'environnement
    Laravel 12, configuration de la base de données SQLite, scaffolding
    d'authentification via Breeze, création des migrations initiales.
    \item \textbf{Sprint 1 (Espace Client --- Réservations)} : développement
    du module de consultation des chambres (galerie, filtrage par dates),
    du formulaire de réservation avec vérification de disponibilité en temps
    réel via l'endpoint AJAX, et de la politique d'annulation à 24 heures.
    \item \textbf{Sprint 2 (Room Service)} : implémentation de la carte de
    menu groupée par catégorie, création des commandes avec calcul automatique
    du total, et machine à états pour le suivi des livraisons.
    \item \textbf{Sprint 3 (Réclamations)} : module de dépôt de réclamations
    avec priorisation automatique, espace d'administration avec assignation
    à un technicien et horodatage de résolution.
    \item \textbf{Sprint 4 (Administration et finalisation)} : gestion des
    réservations côté admin (confirmation/annulation), intégration de
    l'authentification OAuth Google, tests fonctionnels et corrections.
\end{itemize}

\subsection{Planification du projet (diagramme de Gantt)}

La Figure~\ref{fig:gantt} présente la planification globale du projet sous
forme de diagramme de Gantt, répartie sur seize semaines.

\begin{figure}[h]
    \centering
    \begin{ganttchart}[
        hgrid,
        vgrid={*{1}{dotted}},
        x unit=0.72cm,
        y unit title=0.6cm,
        y unit chart=0.55cm,
        title/.style={fill=black!80, draw=none},
        title label font=\color{white}\scriptsize\bfseries,
        bar/.style={fill=black!55, rounded corners=1pt},
        bar label font=\scriptsize,
        group/.style={fill=black!25},
        group label font=\scriptsize\bfseries,
        milestone/.style={fill=black!70, shape=diamond},
        today=8,
        today rule/.style={draw=red, thick, dashed}
    ]{1}{16}
        \gantttitle{Semaines}{16} \\
        \gantttitlelist{1,...,16}{1} \\

        \ganttgroup{Analyse \& Conception}{1}{3} \\
        \ganttbar{Étude des besoins}{1}{2} \\
        \ganttbar{Modélisation UML}{2}{3} \\

        \ganttgroup{Sprint 0 --- Initialisation}{3}{4} \\
        \ganttbar{Env. dev. + migrations}{3}{4} \\

        \ganttgroup{Sprint 1 --- Réservations}{4}{7} \\
        \ganttbar{Chambres + photos}{4}{5} \\
        \ganttbar{Module réservation}{5}{7} \\
        \ganttbar{Vérif. disponibilité (AJAX)}{6}{7} \\

        \ganttgroup{Sprint 2 --- Room Service}{7}{9} \\
        \ganttbar{Carte de menu + commandes}{7}{9} \\

        \ganttgroup{Sprint 3 --- Réclamations}{9}{11} \\
        \ganttbar{Dépôt + priorisation}{9}{11} \\

        \ganttgroup{Sprint 4 --- Administration}{11}{13} \\
        \ganttbar{Espace admin réservations}{11}{12} \\
        \ganttbar{OAuth Google}{12}{13} \\

        \ganttgroup{Tests \& Corrections}{13}{15} \\
        \ganttbar{Tests fonctionnels}{13}{14} \\
        \ganttbar{Corrections et polish}{14}{15} \\

        \ganttgroup{Rédaction du rapport}{1}{16} \\
        \ganttbar{Rédaction continue}{1}{16}
    \end{ganttchart}
    \caption{Diagramme de Gantt du projet Rambla Marrakech}
    \label{fig:gantt}
\end{figure}

\section{Méthodologie et technologies adoptées}

Le développement de la plateforme a suivi une approche itérative : chaque
fonctionnalité a été conçue, implémentée, puis vérifiée de manière
indépendante avant intégration dans l'ensemble du système. Cette méthode
garantit la cohérence des interfaces entre modules et facilite la détection
précoce des anomalies.

Les technologies retenues pour la partie applicative sont les suivantes.
\textbf{Laravel 12} a été choisi comme framework PHP principal en raison de
son architecture MVC mature, de son ORM Eloquent qui simplifie la gestion
des relations polymorphiques, et de son écosystème riche (Breeze pour
l'authentification, Socialite pour OAuth, Tinker pour les tests en console).
\textbf{SQLite} est utilisé comme SGBD en développement pour sa légèreté
et l'absence de configuration serveur. \textbf{Tailwind CSS} assure la
cohérence visuelle via un système de design basé sur des tokens de couleurs
personnalisés. \textbf{Alpine.js} gère les interactions côté client (galerie
d'images, états dynamiques) sans dépendance lourde. \textbf{Flatpickr}
(intégré via CDN) fournit les sélecteurs de dates et d'heures avec
désactivation dynamique des plages déjà réservées. \textbf{Laravel
Socialite} assure l'authentification via OAuth 2.0 avec Google.

Pour la partie réseau, \textbf{Cisco Packet Tracer} a été retenu comme
simulateur réseau, permettant de modéliser la topologie VLAN, les ACL et
la configuration des commutateurs et routeurs de l'hôtel.

\section{Conclusion}

Ce chapitre a permis de poser le cadre du projet en précisant le contexte
hôtelier, la double problématique réseau et applicative, ainsi que les
exigences fonctionnelles et non fonctionnelles retenues. Les technologies
sélectionnées, notamment Laravel 12, Tailwind CSS et Cisco Packet Tracer,
ont été justifiées au regard des besoins identifiés. Les chapitres suivants
détaillent successivement la conception réseau, puis la conception et le
développement de la plateforme web.

% ═════════════════════════════════════════════════════════════════════════
% CHAPITRE 2

\section{Introduction}

xxxx xxxx xxxx xxxx xxxx xxxx xxxx xxxx xxxx xxxx xxxx xxxx xxxx.

\section{Analyse des besoins réseau}

xxxx xxxx xxxx xxxx xxxx xxxx xxxx xxxx xxxx xxxx xxxx xxxx xxxx.

\section{Segmentation VLAN par chambre/étage}

xxxx xxxx xxxx xxxx xxxx xxxx xxxx xxxx xxxx xxxx xxxx xxxx xxxx.

\subsection{Plan d'adressage IP (VLSM)}

xxxx xxxx xxxx xxxx xxxx xxxx xxxx xxxx xxxx xxxx xxxx xxxx xxxx.

\subsection{Sécurité inter-VLAN (ACL)}

xxxx xxxx xxxx xxxx xxxx xxxx xxxx xxxx xxxx xxxx xxxx xxxx xxxx.

\section{Topologie réseau (Packet Tracer)}

xxxx xxxx xxxx xxxx xxxx xxxx xxxx xxxx xxxx xxxx xxxx xxxx xxxx.

\section{Sécurité du Wi-Fi et authentification}

xxxx xxxx xxxx xxxx xxxx xxxx xxxx xxxx xxxx xxxx xxxx xxxx xxxx.

\section{Conclusion}

xxxx xxxx xxxx xxxx xxxx xxxx xxxx xxxx xxxx xxxx xxxx xxxx xxxx.

% ═════════════════════════════════════════════════════════════════════════
% CHAPITRE 3 — CONCEPTION ET DÉVELOPPEMENT DE LA PLATEFORME
% ═════════════════════════════════════════════════════════════════════════

\chapter{Conception et développement de la plateforme e-services}

\section{Introduction}

Ce chapitre présente l'ensemble du travail de conception et de développement
de la plateforme web Rambla Marrakech. Il couvre l'analyse des besoins
fonctionnels, la modélisation du domaine, les choix d'architecture, puis
décrit en détail chacun des trois espaces utilisateurs et les mécanismes de
sécurité mis en \oe uvre.

\section{Analyse et conception}

La phase d'analyse a permis d'identifier quatre profils d'utilisateurs aux
périmètres fonctionnels distincts : le \textbf{client}, hébergé dans une
chambre de l'hôtel, qui consulte les chambres disponibles, effectue des
réservations, passe des commandes de room service et dépose des
réclamations ; le \textbf{serveur}, membre du personnel de restaurant,
dont le rôle se limite au suivi et à la mise à jour des commandes en cours ;
le \textbf{technicien}, chargé du traitement des réclamations techniques
qui lui sont assignées ; et l'\textbf{administrateur}, qui dispose d'une
vision globale sur les réservations et les réclamations, avec la capacité
de les confirmer, les annuler ou les résoudre.

\subsection{Diagramme de cas d'utilisation}

Le diagramme de cas d'utilisation global de la plateforme identifie les
acteurs suivants : \textit{Client}, \textit{Serveur}, \textit{Technicien}
et \textit{Administrateur}. Le client interagit avec les cas d'utilisation
\textit{Consulter les chambres}, \textit{Réserver une chambre ou un service},
\textit{Annuler une réservation}, \textit{Passer une commande de room service}
et \textit{Déposer une réclamation}. Le serveur est associé au cas
\textit{Gérer les commandes}. Le technicien et l'administrateur partagent
les cas \textit{Traiter les réclamations} et \textit{Gérer les réservations}.
L'administrateur possède en outre un accès au tableau de bord de synthèse.

\subsection{Modèle de données (diagramme de classes)}

La base de données de la plateforme comporte dix tables principales, dont
le schéma a été défini au moyen des migrations Laravel. Le Tableau~\ref{tab:tables}
récapitule les entités principales et leurs attributs clés.

\begin{table}[h]
    \centering
    \caption{Entités principales du modèle de données}
    \label{tab:tables}
    \begin{tabular}{|l|p{9.5cm}|}
        \hline
        \textbf{Table} & \textbf{Attributs principaux} \\
        \hline
        \texttt{users} & id, name, email, password, role (client|serveur|technicien|admin), room\_id \\
        \hline
        \texttt{rooms} & id, number, floor, type (simple|double|suite), status (libre|occupee|maintenance), price\_per\_night, wifi\_vlan, surface\_m2, capacity, bed\_type, amenities (JSON), description \\
        \hline
        \texttt{room\_photos} & id, room\_id, url, position \\
        \hline
        \texttt{services} & id, name, description, price, capacity\_per\_slot, is\_active \\
        \hline
        \texttt{reservations} & id, user\_id, reservable\_type, reservable\_id, start\_datetime, end\_datetime, status (en\_attente|confirmee|annulee|terminee), total\_price, notes \\
        \hline
        \texttt{menu\_items} & id, name, description, price, category, is\_available \\
        \hline
        \texttt{orders} & id, user\_id, room\_id, served\_by, status (recue|en\_preparation|en\_livraison|livree|annulee), total\_amount, notes \\
        \hline
        \texttt{order\_items} & id, order\_id, menu\_item\_id, quantity, unit\_price \\
        \hline
        \texttt{complaints} & id, user\_id, room\_id, assigned\_to, type, description, priority (basse|moyenne|urgente), status (ouverte|en\_cours|resolue), resolved\_at \\
        \hline
    \end{tabular}
\end{table}

La relation entre \texttt{reservations} et son entité réservable est
implémentée par un polymorphisme Eloquent (\textit{MorphTo}) : le champ
\texttt{reservable\_type} contient le nom de classe qualifié (\texttt{App\textbackslash{}Models\textbackslash{}Room}
ou \texttt{App\textbackslash{}Models\textbackslash{}Service}), et
\texttt{reservable\_id} l'identifiant correspondant. Ce mécanisme permet
d'unifier la logique de réservation quelle que soit la nature de l'entité
réservée.

% ── Diagramme de cas d'utilisation ──────────────────────────────────────

La Figure~\ref{fig:usecase} présente le diagramme de cas d'utilisation
global de la plateforme, identifiant les quatre acteurs et leurs interactions
avec le système.

\begin{figure}[h]
    \centering
    \begin{tikzpicture}[
        actor/.style={
            text centered, font=\small,
            minimum width=1.8cm
        },
        usecase/.style={
            ellipse, draw, fill=black!8,
            text centered, font=\small,
            minimum width=3.4cm, minimum height=0.65cm,
            inner sep=1pt
        },
        system/.style={
            rectangle, draw, thick,
            inner sep=10pt, rounded corners=4pt
        }
    ]
        % ── Acteurs (gauche) ──
        \node[actor] (client) at (-5.5, 2.8) {
            \tikz\draw[fill=black!15] (0,0) circle (0.28cm);\\[2pt]
            \textbf{Client}
        };
        \node[actor] (serveur) at (-5.5, 0) {
            \tikz\draw[fill=black!15] (0,0) circle (0.28cm);\\[2pt]
            \textbf{Serveur}
        };
        \node[actor] (admin) at (-5.5, -2.8) {
            \tikz\draw[fill=black!15] (0,0) circle (0.28cm);\\[2pt]
            \textbf{Admin / Tech.}
        };

        % ── Cas d'utilisation ──
        \node[usecase] (uc1) at (0,  4.4) {Consulter les chambres};
        \node[usecase] (uc2) at (0,  3.2) {Réserver (chambre/service)};
        \node[usecase] (uc3) at (0,  2.0) {Annuler une réservation};
        \node[usecase] (uc4) at (0,  0.8) {Passer une commande};
        \node[usecase] (uc5) at (0, -0.4) {Déposer une réclamation};
        \node[usecase] (uc6) at (0, -1.6) {Gérer les commandes};
        \node[usecase] (uc7) at (0, -2.8) {Gérer les réservations};
        \node[usecase] (uc8) at (0, -4.0) {Traiter les réclamations};
        \node[usecase] (uc9) at (0, -5.2) {Tableau de bord};

        % ── Frontière système ──
        \begin{scope}[on background layer]
            \node[system, fit=(uc1)(uc2)(uc3)(uc4)(uc5)(uc6)(uc7)(uc8)(uc9),
                  label={[font=\small\bfseries]above:Plateforme Rambla Marrakech}
            ] (sys) {};
        \end{scope}

        % ── Associations Client ──
        \draw (client) -- (uc1);
        \draw (client) -- (uc2);
        \draw (client) -- (uc3);
        \draw (client) -- (uc4);
        \draw (client) -- (uc5);

        % ── Associations Serveur ──
        \draw (serveur) -- (uc6);

        % ── Associations Admin ──
        \draw (admin) -- (uc7);
        \draw (admin) -- (uc8);
        \draw (admin) -- (uc9);
    \end{tikzpicture}
    \caption{Diagramme de cas d'utilisation de la plateforme Rambla Marrakech}
    \label{fig:usecase}
\end{figure}

% ── Diagramme de classes ─────────────────────────────────────────────────

La Figure~\ref{fig:classes} présente le diagramme de classes simplifié
mettant en évidence les entités principales et leurs relations.

\begin{figure}[h]
    \centering
    \begin{tikzpicture}[
        class/.style={
            rectangle, draw, thick, fill=black!6,
            font=\scriptsize,
            minimum width=3.0cm
        },
        header/.style={
            rectangle, draw, thick, fill=black!30,
            text=white, font=\scriptsize\bfseries,
            minimum width=3.0cm, minimum height=0.55cm,
            text centered
        },
        rel/.style={->, >=Stealth, thick},
        comp/.style={->, >=diamond, thick},
        poly/.style={->, >=open diamond, thick, dashed}
    ]
        % ── users ──
        \node[header] (uh) at (0, 8.5) {users};
        \node[class, anchor=north] (u) at (0, 8.18)
            {id, name, email, role\\room\_id (FK)};

        % ── rooms ──
        \node[header] (rh) at (4.8, 8.5) {rooms};
        \node[class, anchor=north] (r) at (4.8, 8.18)
            {id, number, floor, type\\status, price\_per\_night\\wifi\_vlan, amenities (JSON)};

        % ── room\_photos ──
        \node[header] (ph) at (4.8, 5.8) {room\_photos};
        \node[class, anchor=north] (p) at (4.8, 5.48)
            {id, room\_id (FK)\\url, position};

        % ── reservations ──
        \node[header] (resh) at (0, 5.8) {reservations};
        \node[class, anchor=north] (res) at (0, 5.48)
            {id, user\_id (FK)\\reservable\_type\\reservable\_id\\start/end\_datetime\\status, total\_price};

        % ── services ──
        \node[header] (svh) at (9.2, 8.5) {services};
        \node[class, anchor=north] (sv) at (9.2, 8.18)
            {id, name, price\\is\_active};

        % ── orders ──
        \node[header] (oh) at (0, 2.4) {orders};
        \node[class, anchor=north] (o) at (0, 2.08)
            {id, user\_id (FK)\\room\_id (FK), served\_by\\status, total\_amount};

        % ── order\_items ──
        \node[header] (oih) at (4.8, 2.4) {order\_items};
        \node[class, anchor=north] (oi) at (4.8, 2.08)
            {id, order\_id (FK)\\menu\_item\_id (FK)\\quantity, unit\_price};

        % ── complaints ──
        \node[header] (ch) at (0, -0.5) {complaints};
        \node[class, anchor=north] (c) at (0, -0.82)
            {id, user\_id (FK)\\room\_id (FK), assigned\_to\\type, priority, status\\resolved\_at};

        % ── Relations ──
        % users → rooms (belongsTo)
        \draw[rel] (u.east) -- node[above,font=\scriptsize]{occupe} (r.west);
        % rooms ←→ room_photos (hasMany)
        \draw[comp] (r.south) -- node[right,font=\scriptsize]{1..n} (p.north);
        % reservations → users
        \draw[rel] (res.north) -- node[right,font=\scriptsize]{n..1} (u.south);
        % reservations ↔ rooms (morphTo)
        \draw[poly] (res.east) -- node[above,font=\scriptsize]{réservable} (r.south west);
        % reservations ↔ services (morphTo)
        \draw[poly] (res.north east) to[bend left=15] node[above,font=\scriptsize]{réservable} (sv.south west);
        % orders → users
        \draw[rel] (o.north) to[bend right=10] node[left,font=\scriptsize]{n..1} (u.south);
        % orders → order_items (hasMany)
        \draw[comp] (o.east) -- node[above,font=\scriptsize]{1..n} (oi.west);
        % complaints → users
        \draw[rel] (c.north) to[bend left=10] node[left,font=\scriptsize]{n..1}(o.south);
    \end{tikzpicture}
    \caption{Diagramme de classes simplifié de la plateforme Rambla Marrakech}
    \label{fig:classes}
\end{figure}

\section{Choix technologiques}

L'architecture logicielle de la plateforme repose sur le patron
Modèle-Vue-Contrôleur (MVC) implémenté par le framework \textbf{Laravel 12}.
Ce choix s'explique par plusieurs raisons concrètes : la gestion native des
migrations de base de données garantit une évolution contrôlée du schéma ;
l'ORM Eloquent offre une abstraction cohérente des relations entre entités,
y compris les relations polymorphiques nécessaires au module de réservation ;
le moteur de templates Blade simplifie la séparation entre logique métier et
présentation ; enfin, le système de middleware permet d'appliquer les règles
de contrôle d'accès de manière centralisée et réutilisable.

\textbf{SQLite} est retenu comme système de gestion de base de données en
environnement de développement. Sa configuration ne nécessite aucun serveur
dédié, ce qui simplifie le déploiement initial et les phases de test.

\textbf{Tailwind CSS} est utilisé pour la construction de l'interface
utilisateur. Le projet définit un système de design propre, basé sur des
tokens de couleurs nommés (\texttt{ink}, \texttt{parchment}, \texttt{brass},
\texttt{coral}, \texttt{sage}, \texttt{stone}), des classes utilitaires
personnalisées (\texttt{hotel-panel}, \texttt{key-card}, \texttt{status-badge})
et deux polices de caractères : Fraunces (serif, pour les titres) et
Public Sans (sans-serif, pour le corps de texte).

\textbf{Alpine.js} est intégré pour les comportements interactifs côté
client ne nécessitant pas de rechargement de page : galerie photographique
avec transitions, menus déroulants, états d'affichage conditionnels.

\textbf{Flatpickr} (bibliothèque légère chargée via CDN) est utilisé pour
remplacer les champs \texttt{datetime-local} natifs par des sélecteurs de
dates visuels. Les plages de dates déjà réservées sont transmises au
composant en JSON par le contrôleur et désactivées dynamiquement dans le
calendrier au moment de la sélection de la chambre.

\textbf{Laravel Socialite} assure l'intégration OAuth 2.0 avec Google,
permettant aux clients de s'authentifier sans créer de mot de passe local.
L'utilisateur est créé automatiquement lors du premier accès avec le rôle
\texttt{client} et un mot de passe aléatoire non utilisé.

\section{Les trois espaces de l'application}

\subsection{Espace Client}

L'espace client regroupe l'ensemble des fonctionnalités accessibles à un
utilisateur authentifié avec le rôle \texttt{client}. Il est sécurisé par
le groupe de routes \texttt{middleware(['auth', 'role:client'])} et préfixé
par \texttt{/client}. Le tableau de bord (\texttt{GET /client/dashboard})
présente les commandes actives et les réclamations non résolues de
l'utilisateur.

\textbf{Consultation et réservation des chambres.} Le module de consultation
(\texttt{RoomController@index}) liste les chambres dont le statut est
\texttt{libre}, ordonnées par type puis par numéro. Un filtre de dates
(\texttt{?start=\&end=}) active le scope \texttt{availableBetween()} défini
dans le modèle \texttt{Room}, qui exclut via \texttt{whereDoesntHave}
toutes les chambres ayant une réservation au statut \texttt{en\_attente}
ou \texttt{confirmee} chevauchant la période demandée. La fiche détaillée
(\texttt{RoomController@show}) expose les caractéristiques complètes de la
chambre : surface en m², capacité, type de lit, liste des équipements
(\textit{amenities}) stockée en colonne JSON, description longue et galerie
de photos ordonnées par la colonne \texttt{position} de la table
\texttt{room\_photos}.

Le formulaire de réservation (\texttt{ReservationController@create})
injecte en JSON deux structures : \texttt{\$roomsJson} contenant les
métadonnées de chaque chambre (numéro, photo de couverture, caractéristiques)
pour la prévisualisation dynamique côté client, et \texttt{\$roomsBookedRanges}
contenant les plages de dates déjà réservées par chambre pour les transmettre
à Flatpickr et y désactiver visuellement les dates indisponibles.

Avant la création effective d'une réservation, \texttt{ReservationController@store}
appelle la méthode \texttt{Room::isAvailable(\$start, \$end)} qui construit
une seule requête combinant les deux cas de chevauchement --- réservation
avec date de fin explicite et réservation sans date de fin (nuit implicite)
--- et applique une négation unique sur le résultat de \texttt{exists()}.
Si un chevauchement est détecté, la soumission est rejetée avec un message
d'erreur explicite. En l'absence de conflit, le prix total est calculé en
multipliant le tarif par nuit par le nombre de nuits, puis la réservation
est créée au statut \texttt{en\_attente}. Un endpoint AJAX
(\texttt{GET /client/reservations/check-availability}) permet également
une vérification en temps réel lors de la saisie des dates, avec retour
JSON \texttt{\{available: true|false\}} et feedback visuel immédiat dans
le formulaire.

La politique d'annulation impose un délai minimal de 24 heures avant
l'arrivée, défini par la constante \texttt{Reservation::CANCELLATION\_DEADLINE\_HOURS}.
La méthode \texttt{isCancellableByClient()} vérifie conjointement le statut
actif de la réservation et ce délai ; le contrôleur s'assure en outre que
la réservation appartient bien à l'utilisateur authentifié avant d'autoriser
l'opération.

\textbf{Room service.} Le module de room service (\texttt{RoomServiceController})
est conditionné à l'assignation préalable du client à une chambre
(\texttt{user.room\_id}). La carte de menu est chargée depuis le modèle
\texttt{MenuItem}, filtrée par le scope \texttt{disponibles()} et regroupée
par catégorie. À la soumission, le contrôleur crée un enregistrement
\texttt{Order} au statut \texttt{recue}, puis un \texttt{OrderItem} par
article sélectionné avec sa quantité et son prix unitaire au moment de la
commande. La méthode \texttt{Order::recalculateTotal()} recalcule le montant
total par agrégation SQL (\texttt{SUM(quantity * unit\_price)}).

\textbf{Réclamations.} La soumission d'une réclamation (\texttt{ComplaintController@store})
requiert que le client soit assigné à une chambre. Le type est validé contre
l'énumération \texttt{Complaint::types()} qui comprend : climatisation,
eau\_chaude, télévision, wifi, plomberie, électricité, autre. La priorité
est déterminée automatiquement : les types figurant dans
\texttt{Complaint::typesUrgentsParDefaut()} (climatisation, eau\_chaude,
plomberie, électricité) reçoivent le niveau \texttt{urgente} ; les autres
reçoivent \texttt{moyenne}. La réclamation est créée au statut
\texttt{ouverte}.

\subsection{Espace Serveur / Restaurant}

L'espace serveur est sécurisé par le middleware \texttt{role:serveur} et
préfixé par \texttt{/staff}. Il expose deux routes uniquement :
\texttt{GET /staff/orders} et \texttt{PATCH /staff/orders/\{order\}}.

Le contrôleur \texttt{Staff\textbackslash{}OrderController} charge en index
toutes les commandes dont le statut est différent de \texttt{livree} ou
\texttt{annulee}, avec chargement eager des relations \texttt{user}, 
\texttt{room} et \texttt{items.menuItem}. La mise à jour d'une commande
applique une machine à états séquentielle définie par le tableau de
constante \texttt{NEXT\_STATUS} :

\begin{center}
\texttt{recue} $\rightarrow$ \texttt{en\_preparation} $\rightarrow$
\texttt{en\_livraison} $\rightarrow$ \texttt{livree}
\end{center}

Toute tentative de transition vers un état non prévu par cette séquence
est rejetée avec un message d'erreur de validation. À chaque mise à jour,
l'identifiant du serveur traitant la commande est enregistré dans la colonne
\texttt{served\_by}. Cette conception garantit l'intégrité du cycle de vie
d'une commande et facilite la traçabilité des opérations.

\subsection{Espace Administration / Technique}

L'espace d'administration est accessible aux utilisateurs portant le rôle
\texttt{technicien} ou \texttt{admin}, conformément au middleware
\texttt{role:technicien,admin}. Il est préfixé par \texttt{/admin}.

Le tableau de bord (\texttt{GET /admin/dashboard}) présente le nombre de
réclamations au statut \texttt{urgente} qui sont encore ouvertes ou en
cours, calculé par la combinaison des scopes \texttt{Complaint::urgentes()}
et \texttt{whereIn('status', ['ouverte', 'en\_cours'])}.

\textbf{Gestion des réclamations.} Le contrôleur
\texttt{Admin\textbackslash{}ComplaintController@index} charge toutes les
réclamations avec leurs relations (\texttt{user}, \texttt{room},
\texttt{assignedTo}), triées par priorité décroissante (urgente en premier)
puis par date de création décroissante, via une clause \texttt{orderByRaw}
avec une expression \texttt{CASE}. L'action \texttt{update} permet
d'assigner la réclamation à un technicien et de faire évoluer son statut
vers \texttt{en\_cours} ou \texttt{resolue}. La transition vers
\texttt{resolue} déclenche la méthode \texttt{Complaint::marquerResolue()}
qui enregistre l'horodatage dans la colonne \texttt{resolved\_at}.

\textbf{Gestion des réservations.} Le contrôleur
\texttt{Admin\textbackslash{}ReservationController@index} liste toutes
les réservations avec leurs relations (\texttt{user}, \texttt{reservable}),
les réservations \texttt{en\_attente} remontant en première position grâce
à un ordre composite. L'action \texttt{update} valide que le nouveau statut
est strictement \texttt{confirmee} ou \texttt{annulee} avant d'appliquer
la modification.

\section{Sécurité et gestion des rôles}

La sécurisation de la plateforme repose sur plusieurs mécanismes
complémentaires.

\textbf{Authentification.} Le scaffolding d'authentification est fourni par
Laravel Breeze, qui génère les routes, contrôleurs et vues pour
l'inscription, la connexion par formulaire, la réinitialisation de mot de
passe et la vérification d'adresse email. La connexion via OAuth Google est
gérée par \texttt{GoogleAuthController} : après redirection vers Google et
retour du callback, l'utilisateur est recherché par son adresse email via
\texttt{User::firstOrCreate()} ; s'il n'existe pas, il est créé avec le
rôle \texttt{client} et un mot de passe aléatoire hashé ; il est ensuite
connecté via \texttt{Auth::login()} et redirigé selon son rôle.

\textbf{Contrôle d'accès par rôle.} Le middleware \texttt{CheckRole}
reçoit un ou plusieurs rôles en paramètre (\texttt{role:client},
\texttt{role:technicien,admin}, etc.) et renvoie une réponse HTTP 403 si
le rôle de l'utilisateur authentifié ne figure pas dans la liste autorisée.
Ce middleware est appliqué au niveau du groupe de routes, ce qui évite toute
duplication de la logique d'autorisation dans les contrôleurs. La colonne
\texttt{role} de la table \texttt{users} est une chaîne de caractères non
chiffrée, directement interrogeable.

\textbf{Protection CSRF.} Chaque formulaire POST, PATCH ou DELETE inclut
un jeton \texttt{@csrf} généré par Blade. Les routes de type API légère
(\texttt{check-availability}) acceptent les requêtes \texttt{XMLHttpRequest}
sans CSRF en raison de leur nature GET, mais renvoient uniquement des
données en lecture.

\textbf{Validation des entrées.} Chaque contrôleur utilise la méthode
\texttt{\$request->validate()} de Laravel avec des règles explicites :
types d'enum validés via \texttt{Rule::in()}, existence en base vérifiée
par la règle \texttt{exists:}, longueurs maximales imposées sur les champs
textuels. Les erreurs de validation sont renvoyées automatiquement vers le
formulaire avec les valeurs saisies préservées (\texttt{old()}).

\textbf{Isolation des données.} La méthode \texttt{cancel} du contrôleur
de réservation vérifie explicitement que \texttt{\$reservation->user\_id ===
auth()->id()} avant d'autoriser l'opération, empêchant ainsi qu'un client
puisse annuler la réservation d'un autre utilisateur.

\section{Conclusion}

Ce chapitre a présenté l'architecture complète de la plateforme web Rambla
Marrakech. Le modèle de données repose sur neuf tables interconnectées,
dont une relation polymorphique qui unifie la réservation de chambres et
de services. Les trois espaces utilisateurs --- client, serveur et
administration --- sont strictement délimités par des groupes de routes
protégés par middleware. La logique métier critique, notamment la
vérification des chevauchements de réservations et la priorisation
automatique des réclamations, est encapsulée dans les modèles Eloquent
sous forme de méthodes et de scopes réutilisables. Le chapitre suivant
présente les tests réalisés sur ces fonctionnalités et les résultats
obtenus. — TESTS, RÉSULTATS ET PERSPECTIVES
% ═════════════════════════════════════════════════════════════════════════

\chapter{Tests, résultats et perspectives}

\section{Introduction}

Ce chapitre présente la démarche de validation adoptée pour vérifier le
comportement de la plateforme. Il décrit les scénarios de tests fonctionnels
exécutés sur les modules principaux, expose les résultats obtenus, identifie
les difficultés rencontrées durant le développement et propose des
perspectives d'amélioration pour les versions ultérieures.

\section{Scénarios de tests fonctionnels}

Les tests ont été conduits manuellement en environnement de développement
à l'aide de l'outil \texttt{php artisan tinker} pour les vérifications de
logique métier et du navigateur pour les flux d'interface. Le Tableau~\ref{tab:tests}
récapitule les principaux scénarios testés.

\begin{table}[h]
    \centering
    \caption{Scénarios de tests fonctionnels}
    \label{tab:tests}
    \begin{tabular}{|p{5.5cm}|p{4cm}|p{4cm}|}
        \hline
        \textbf{Scénario} & \textbf{Entrées / conditions} & \textbf{Résultat attendu} \\
        \hline
        Connexion par formulaire & Identifiants valides & Redirection vers le tableau de bord du rôle correspondant \\
        \hline
        Connexion OAuth Google & Compte Google valide & Création ou récupération de l'utilisateur, redirection client \\
        \hline
        Accès à un espace non autorisé & URL \texttt{/admin} avec rôle \texttt{client} & Réponse HTTP 403 \\
        \hline
        Réservation sans conflit de dates & Chambre libre, dates non chevauchantes & Réservation créée au statut \texttt{en\_attente} \\
        \hline
        Réservation avec conflit de dates & Chambre déjà réservée sur la période & Rejet avec message d'erreur explicite \\
        \hline
        Vérification AJAX de disponibilité & \texttt{GET /check-availability} avec dates valides & JSON \texttt{\{available: true\}} ou \texttt{\{available: false\}} \\
        \hline
        Annulation dans le délai & Réservation active, arrivée à plus de 24 h & Statut mis à jour vers \texttt{annulee} \\
        \hline
        Annulation hors délai & Réservation active, arrivée à moins de 24 h & Rejet avec message invitant à contacter la réception \\
        \hline
        Tentative d'annulation d'autrui & \texttt{user\_id} différent de l'authentifié & Réponse HTTP 403 \\
        \hline
        Commande room service & Sélection de plusieurs articles & Création \texttt{Order} + \texttt{OrderItem}, total recalculé \\
        \hline
        Avancement de commande hors séquence & Transition \texttt{recue} $\rightarrow$ \texttt{livree} & Rejet : transition invalide \\
        \hline
        Dépôt de réclamation type urgent & Type \texttt{climatisation} & Priorité automatique \texttt{urgente} \\
        \hline
        Dépôt de réclamation type standard & Type \texttt{wifi} & Priorité automatique \texttt{moyenne} \\
        \hline
        Résolution d'une réclamation & Statut passé à \texttt{resolue} & \texttt{resolved\_at} horodaté à l'instant courant \\
        \hline
        Confirmation de réservation (admin) & Statut \texttt{en\_attente} $\rightarrow$ \texttt{confirmee} & Mise à jour persistée, flash de succès \\
        \hline
    \end{tabular}
\end{table}

\section{Résultats obtenus}

L'ensemble des scénarios décrits dans le Tableau~\ref{tab:tests} ont
produit les résultats conformes aux spécifications. Les points notables
sont les suivants.

La logique de détection des chevauchements de réservations a été validée
sur huit cas distincts, couvrant les chevauchements partiels en début et
en fin de période, les englobements complets, les périodes non chevauchantes
avant et après une réservation existante, et le cas des réservations dont
le statut est \texttt{annulee} (qui ne doivent pas bloquer les nouvelles
demandes). Tous ces cas ont retourné le résultat attendu.

Le mécanisme de machine à états pour les commandes de room service a
correctement rejeté toutes les transitions non contiguës dans la séquence
définie. L'enregistrement de l'identifiant du serveur traitant (\texttt{served\_by})
a été vérifié à chaque transition.

La priorisation automatique des réclamations a assigné le niveau
\texttt{urgente} aux types climatisation, eau\_chaude, plomberie et
électricité, et le niveau \texttt{moyenne} aux types télévision, wifi et
autre, conformément à la méthode \texttt{typesUrgentsParDefaut()} du
modèle \texttt{Complaint}.

La vérification en temps réel de la disponibilité via l'endpoint AJAX a
fonctionné correctement dans le formulaire de réservation, avec désactivation
du bouton de soumission en cas d'indisponibilité détectée avant même la
soumission du formulaire.

Le filtrage de disponibilité sur la page de liste des chambres
(\texttt{/client/rooms?start=\&end=}) a retourné uniquement les chambres
sans conflit sur la période saisie, grâce au scope
\texttt{Room::availableBetween()}.

\section{Difficultés rencontrées}

Plusieurs difficultés techniques ont nécessité une analyse approfondie avant
d'être résolues.

\textbf{Logique booléenne de détection des chevauchements.} La première
implémentation de la méthode \texttt{isAvailable()} combinait deux appels
à \texttt{exists()} distincts — un pour les réservations avec date de fin
explicite et un pour les réservations sans date de fin — reliés par une
expression \texttt{!\$A || !\$B}. Cette structure est mathématiquement
équivalente à \texttt{!(\$A \&\& \$B)}, ce qui rend la condition vraie dès
qu'un des deux sous-ensembles est vide, produisant de faux positifs de
disponibilité. La correction a consisté à regrouper les deux cas dans un
seul \texttt{where()} avec \texttt{orWhere()} imbriqué, et à appliquer la
négation une seule fois sur le \texttt{exists()} global.

\textbf{Structure de la vue de réservation.} Lors de l'ajout de la
bibliothèque Flatpickr, la balise \texttt{<x-app-layout>} a été
accidentellement omise de la vue \texttt{create.blade.php}, entraînant un
rendu sans layout (absence du \texttt{<head>}, de la navigation et des
feuilles de style). Cette anomalie a été corrigée en rétablissant la
balise ouvrante après le premier bloc \texttt{@push('head-assets')} et la
balise fermante en fin de fichier.

\textbf{Gestion des paramètres de route.} La route AJAX
\texttt{/client/reservations/check-availability} devait être déclarée avant
la route nommée \texttt{/client/reservations/create} afin d'éviter que le
routeur de Laravel ne l'interprète comme un paramètre de segment d'URL.

\textbf{Compatibilité du sélecteur de dates sur différents navigateurs.}
Les champs \texttt{datetime-local} natifs produisent des rendus visuels
hétérogènes selon le navigateur. Leur remplacement par Flatpickr a unifié
l'expérience utilisateur, mais a nécessité l'ajout de la propriété CSS
\texttt{[color-scheme:light]} pour forcer le thème clair sur les navigateurs
appliquant un thème sombre système.

\section{Perspectives d'amélioration}

Plusieurs axes d'évolution peuvent être envisagés pour enrichir la
plateforme dans ses versions ultérieures.

\textbf{Notifications en temps réel.} L'intégration de WebSockets via
Laravel Reverb ou Pusher permettrait d'informer le client en temps réel
de la confirmation ou de l'annulation de sa réservation par
l'administration, sans recours au rechargement de page.

\textbf{Paiement en ligne.} L'ajout d'une passerelle de paiement (Stripe,
CMI ou PayPal) au module de réservation permettrait de collecter un acompte
ou le paiement total au moment de la confirmation, renforçant la fiabilité
des réservations.

\textbf{Tableau de bord analytique.} Un module de statistiques exposant
le taux d'occupation par type de chambre, la distribution des réclamations
par catégorie et l'évolution du chiffre d'affaires mensuel apporterait une
valeur décisionnelle à l'espace administrateur.

\textbf{Application mobile.} La transformation de la plateforme en une
Progressive Web App (PWA) ou l'exposition d'une API REST consommée par une
application native permettrait aux clients d'accéder aux services depuis
leur téléphone sans installation préalable.

\textbf{Gestion des disponibilités de services.} Le modèle \texttt{Service}
dispose d'un champ \texttt{capacity\_per\_slot} non encore exploité côté
réservation. Son activation permettrait de limiter le nombre de réservations
simultanées sur un même créneau pour un service donné.

\textbf{Internationalisation.} L'ajout du support multilingue (arabe,
anglais, français) via le système de localisation de Laravel rendrait la
plateforme accessible à un public international.

\section{Conclusion}

Ce chapitre a validé le comportement de la plateforme sur les quinze
scénarios fonctionnels couvrant les modules de réservation, de room service,
de réclamations et de contrôle d'accès. Les résultats obtenus sont conformes
aux spécifications définies au chapitre premier. Les difficultés rencontrées,
principalement d'ordre logique et structurel, ont conduit à des corrections
documentées qui renforcent la robustesse du code. Les perspectives
identifiées offrent des axes concrets d'évolution vers une plateforme plus
complète et mieux intégrée dans l'écosystème numérique hôtelier.
% ─────────────────────────────────────────────────────────────────────────

\chapter*{Conclusion générale et perspectives}
\addcontentsline{toc}{chapter}{Conclusion générale et perspectives}

Ce projet a conduit à la réalisation de deux livrables complémentaires
répondant aux problématiques identifiées dans le cadre de l'hôtel Rambla
Marrakech. Sur le plan réseau, une architecture VLAN par chambre a été
conçue et simulée sous Cisco Packet Tracer, avec segmentation des flux,
filtrage inter-VLAN par ACL et sécurisation du réseau Wi-Fi. Sur le plan
applicatif, une plateforme web développée sous Laravel 12 couvre l'ensemble
du cycle de vie des services hôteliers : réservation de chambres et de
prestations avec détection des conflits de disponibilité, room service avec
machine à états séquentielle, réclamations techniques avec priorisation
automatique et workflow de résolution, et authentification multi-méthodes
incluant OAuth Google.

L'architecture MVC de Laravel, combinée à l'utilisation rigoureuse des
scopes Eloquent, des relations polymorphiques et du middleware de contrôle
d'accès, a permis de construire une base de code structurée et extensible.
La cohérence de l'interface utilisateur a été assurée par un système de
design custom fondé sur des tokens de couleurs et des classes utilitaires
Tailwind CSS réutilisables sur l'ensemble des vues.

Les perspectives d'amélioration identifiées --- notifications en temps réel,
paiement en ligne, tableau de bord analytique, application mobile et
internationalisation --- constituent des extensions naturelles du périmètre
actuel, réalisables sans remise en cause de l'architecture existante.

% ─────────────────────────────────────────────────────────────────────────
% RÉFÉRENCES / BIBLIOGRAPHIE
% ─────────────────────────────────────────────────────────────────────────

\begin{thebibliography}{99}
    \addcontentsline{toc}{chapter}{Références}

    \bibitem{laravel} Laravel Documentation. Disponible à l'adresse :
    \url{https://laravel.com/docs} (consulté le XX 2026).

    \bibitem{ciscopt} Cisco Packet Tracer. Documentation officielle.
    Disponible à l'adresse : \url{https://www.netacad.com/} (consulté le
    XX 2026).

    \bibitem{tailwind} Tailwind CSS Documentation. Disponible à l'adresse :
    \url{https://tailwindcss.com/docs} (consulté le XX 2026).

\end{thebibliography}

% ─────────────────────────────────────────────────────────────────────────
% ANNEXES
% ─────────────────────────────────────────────────────────────────────────

\appendix

\chapter*{Annexes}
\addcontentsline{toc}{chapter}{Annexes}

\section*{Annexe A : Migrations et schéma de base de données}

La plateforme compte quatorze fichiers de migration Laravel exécutés dans
l'ordre chronologique défini par leur horodatage de nommage. Le
Tableau~\ref{tab:migrations} liste ces migrations dans leur ordre
d'exécution.

\begin{table}[h]
    \centering
    \caption{Liste des migrations de la base de données}
    \label{tab:migrations}
    \begin{tabular}{|l|l|}
        \hline
        \textbf{Fichier de migration} & \textbf{Opération} \\
        \hline
        \texttt{0001\_01\_01\_create\_users\_table} & Création de la table \texttt{users} \\
        \hline
        \texttt{0001\_01\_01\_create\_cache\_table} & Création des tables de cache \\
        \hline
        \texttt{0001\_01\_01\_create\_jobs\_table} & Création de la table de jobs \\
        \hline
        \texttt{2026\_07\_08\_create\_rooms\_table} & Création de \texttt{rooms}, ajout de \texttt{room\_id} sur \texttt{users} \\
        \hline
        \texttt{2026\_07\_08\_create\_services\_table} & Création de \texttt{services} \\
        \hline
        \texttt{2026\_07\_08\_create\_reservations\_table} & Création de \texttt{reservations} avec index polymorphe \\
        \hline
        \texttt{2026\_07\_08\_create\_menu\_items\_table} & Création de \texttt{menu\_items} \\
        \hline
        \texttt{2026\_07\_08\_create\_orders\_table} & Création de \texttt{orders} \\
        \hline
        \texttt{2026\_07\_08\_create\_order\_items\_table} & Création de \texttt{order\_items} \\
        \hline
        \texttt{2026\_07\_08\_create\_complaints\_table} & Création de \texttt{complaints} \\
        \hline
        \texttt{2026\_07\_20\_add\_characteristics\_to\_rooms} & Ajout de \texttt{surface\_m2}, \texttt{capacity}, \texttt{bed\_type}, \texttt{amenities}, \texttt{description} \\
        \hline
        \texttt{2026\_07\_20\_create\_room\_photos\_table} & Création de \texttt{room\_photos} \\
        \hline
    \end{tabular}
\end{table}

Les relations entre tables suivent le schéma ci-après. La table
\texttt{users} référence \texttt{rooms} via la clé étrangère
\texttt{room\_id} (nullable, \texttt{nullOnDelete}). La table
\texttt{reservations} utilise une clé polymorphe composée de
\texttt{reservable\_type} et \texttt{reservable\_id}, indexée
conjointement, permettant de référencer indifféremment une
\texttt{Room} ou un \texttt{Service}. La table \texttt{complaints}
référence \texttt{users} deux fois : via \texttt{user\_id} (auteur)
et \texttt{assigned\_to} (technicien assigné, nullable). La table
\texttt{orders} référence \texttt{users} via \texttt{user\_id} (client)
et \texttt{served\_by} (serveur, nullable). La table \texttt{order\_items}
référence \texttt{orders} et \texttt{menu\_items} avec suppression en
cascade.

\section*{Annexe B : Plan d'adressage IP détaillé}

Cette annexe présentera le plan d'adressage IP complet de l'hôtel Rambla
Marrakech une fois les détails de la topologie Packet Tracer fournis. Elle
comprendra les éléments suivants : le découpage en sous-réseaux par
technique VLSM, le tableau des VLANs avec leur identifiant numérique, leur
nom, leur plage d'adresses, le masque de sous-réseau, la passerelle par
défaut et la capacité en hôtes ; le détail des ACL inter-VLAN définissant
les flux autorisés et refusés entre segments ; et la configuration des
interfaces du routeur inter-VLAN et des commutateurs gérés.

\end{document}
